\chapter{Three-Dimensional Chern--Simons--Matter Theories}
\label{ch:susy-three-dimensional-chern-simons-matter}

\ConventionsMixed

In three dimensions, supersymmetric gauge theories combine topological gauge
terms, parity anomalies, monopole operators, and matter dynamics.  This chapter
fixes the basic \(3D\) \(\mathcal N=2\) Chern--Simons--matter datum before
later protected-sector and localization chapters use it.

\section{Three-Dimensional Supersymmetry}

Let \(\eta_{\mu\nu}=\operatorname{diag}(-,+,+)\).  A minimal Majorana spinor
has two real components.  The \(\mathcal N=2\) algebra may be written with
complex supercharges \(Q_\alpha,\bar Q_\alpha\):
\[
  \{Q_\alpha,\bar Q_\beta\}
  =
  2(\gamma^\mu)_{\alpha\beta}P_\mu
  +2\ii\epsilon_{\alpha\beta}Z,
\]
where \(Z\) is a possible central charge.  Particle multiplets are
representations of this Hilbert-space algebra.  Superfields are off-shell
coordinate packages for local variables; the two notions are connected only
after dynamics and quantization have been specified.

\section{Chern--Simons Datum}

Let \(G\) be compact and let \(P\to M_3\) be a principal \(G\)-bundle.  A
Chern--Simons action is
\[
  S_{\mathrm{CS}}[A]
  =
  \frac{k}{4\pi}
  \int_{M_3}\operatorname{Tr}
  \left(A\wedge dA+\frac23A\wedge A\wedge A\right).
\]
The trace convention and the allowed levels \(k\) are tied to the lattice of
large gauge transformations.  Gauge invariance of
\(\exp(\ii S_{\mathrm{CS}})\) requires the change of \(S_{\mathrm{CS}}\) under
a large gauge transformation to lie in \(2\pi\mathbb Z\).

Varying the action gives
\[
  \delta S_{\mathrm{CS}}
  =
  \frac{k}{2\pi}\int_{M_3}\operatorname{Tr}(\delta A\wedge F_A)
  +
  \frac{k}{4\pi}\int_{\partial M_3}\operatorname{Tr}(A\wedge\delta A).
\]
On a closed manifold, the field equation of pure Chern--Simons theory is
\[
  F_A=0.
\]
On a manifold with boundary, a boundary condition or boundary degrees of
freedom are part of the gauge-theory datum.

\section{Supersymmetric Vector Multiplet}

A \(3D\) \(\mathcal N=2\) vector multiplet contains
\[
  (A_\mu,\sigma,\lambda,\bar\lambda,D),
\]
where \(\sigma\) is a real adjoint scalar and \(D\) is an auxiliary field.
The supersymmetric Chern--Simons action contains
\[
  S_{\mathrm{CS}}^{\mathcal N=2}
  =
  \frac{k}{4\pi}\int
  \operatorname{Tr}
  \left(
    A\wedge dA+\frac23A^3
    -\bar\lambda\lambda
    +2D\sigma
  \right),
\]
with the displayed signs defining the convention of this chapter.  The
auxiliary equation couples \(D\) to matter moment maps and to \(\sigma\),
thereby modifying the vacuum equations compared with pure Yang--Mills matter
systems.

\section{Matter And Moment Maps}

Let chiral multiplets \(\Phi_i\) take values in a unitary representation
\(R\) of \(G\).  Their scalar fields are \(\phi_i\).  With real mass
parameters \(m_i\), the bosonic potential contains
\[
  \sum_i\left|(\sigma+m_i)\phi_i\right|^2
  +
  \frac12\left(\mu(\phi)+\frac{k}{2\pi}\sigma-\zeta\right)^2
  +
  |dW(\phi)|^2,
\]
in an abelian normalization where \(\mu\) is the matter moment map and
\(\zeta\) is the FI coordinate.  Supersymmetric classical vacua solve
\[
  dW=0,\qquad
  (\sigma+m_i)\phi_i=0,\qquad
  \mu(\phi)+\frac{k}{2\pi}\sigma-\zeta=0,
\]
modulo gauge transformations.

\section{Parity Anomaly}

A massive three-dimensional Dirac fermion in representation \(R\) shifts the
effective Chern--Simons level by
\[
  \Delta k=\frac12\,\operatorname{sign}(m)\,T(R),
\]
where
\[
  \operatorname{Tr}_R(t^at^b)=T(R)\delta^{ab}
\]
in the trace convention used for the gauge action.  Therefore a massless
fermion coupled to a gauge field requires a choice of regulator that fixes the
integer or half-integer effective level.  This choice is part of the
definition of the continuum theory and of its time-reversal or parity
properties.

\section{Monopole Operators}

Let \(G\) have maximal torus \(T\).  A monopole operator at the origin is
defined by imposing a magnetic singularity on a small sphere:
\[
  \frac1{2\pi}\int_{S^2}F=m\in\operatorname{Hom}(U(1),T)/W_G .
\]
The local operator is gauge invariant only after its electric charges induced
by Chern--Simons terms and fermion zero modes are accounted for.  In an
abelian theory with level \(k\), Gauss law gives electric charge
\[
  q_{\mathrm{el}}=k\,m
\]
for a bare monopole of magnetic charge \(m\).  Gauge-invariant monopole
operators may require dressing by charged matter fields.

\section{Three-Dimensional Dynamics Ledger}

A complete \(3D\) supersymmetric Chern--Simons--matter example must specify:
\[
\begin{gathered}
  G,\quad k,\quad R,\quad W,\quad \zeta,\quad m_i,\\
  \text{spin structure and framing data},\quad
  \text{parity-anomaly regulator},\\
  \text{monopole-operator dressing rules},\quad
  \text{boundary and global-form choices}.
\end{gathered}
\]
Duality proposals compare these full data, including contact terms and
background Chern--Simons counterterms.

\begin{openproblem}[Three-dimensional supersymmetric continuum schemes]
\label{op:three-dimensional-susy-cs-matter-continuum}
Construct \(3D\) supersymmetric Chern--Simons--matter theories with a
regulator preserving the stated gauge and supersymmetry Ward identities,
including parity-anomaly choices, monopole-operator definitions, boundary
conditions, and contact-term coordinates.
\end{openproblem}
